\subsection{What one density adds to one support}
\label{sec:v27-density-support}

The family of support thresholds determines the action, but the support
at one offset need not do so.  The following functional comparison
isolates the additional information used by the four-density identity.

\begin{proposition}[A support-preserving action change]
\label{prop:v27-support-gauge}
On $[-1,1]$ let $d=1$, $S(u)=u^2$, and, for $0<\epsilon\le1/10$, set
\[
 \phi_\epsilon(s)=s+\epsilon s(1-s)(s-1/2),\qquad
 \widetilde S(u)=\phi_\epsilon(S(u)).
\]
Both actions vanish to second order at zero, are strictly convex, and
are different.  Nevertheless
\begin{equation}\label{eq:v27-equal-support}
 S(u)+S(v)<1\quad\Longleftrightarrow\quad
                    \widetilde S(u)+\widetilde S(v)<1.
\end{equation}
Their endpoint densities with any positive separable amplitudes cannot
have the same four-density invariant on an interior neighborhood of zero.
\end{proposition}

\begin{proof}
For $0\le s\le1$,
\[
 \phi_\epsilon'(s)=1+\epsilon(-3s^2+3s-1/2)\ge1-\epsilon/2>0,
 \qquad \phi_\epsilon(1-s)=1-\phi_\epsilon(s).
\]
Consequently, for $s,t\in[0,1]$, $s+t<1$ is equivalent to
$\phi_\epsilon(t)<\phi_\epsilon(1-s)=1-\phi_\epsilon(s)$.
This proves \eqref{eq:v27-equal-support}.  Also
\[
 \widetilde S''(u)
 =2\{\phi_\epsilon'(s)+2s\phi_\epsilon''(s)\}_{s=u^2},\qquad
 \phi_\epsilon'(s)+2s\phi_\epsilon''(s)
 =1+\epsilon(-15s^2+9s-1/2)\ge1-13\epsilon/2>0.
\]
The zero value and derivative at the origin follow by composition.
For $0<S(u)<1/2$ the correction is nonzero, so the actions differ on
every neighborhood of zero.  On a sufficiently small interior square,
\eqref{eq:v26-rank-one-defect} and the fixed-anchor inverse recover $S$
from its four-density invariant, irrespective of the positive amplitude.
Equality of the invariants would therefore make the two actions equal,
a contradiction.
\end{proof}

This example belongs to the functional action--density class.  It is not
an assertion that both actions have been realized by physical billiard
tables satisfying the global hypotheses.  The physical theorem uses the
relative-law construction for that realization.  The comparison shows
precisely why retaining density, rather than just one support, enters
Theorem~\ref{thm:v26-density-inverse}.
